% Generated draft: requires mathematical and source review.
\documentclass[11pt]{article}
\usepackage[margin=25mm]{geometry}
\usepackage{fontspec}
\setmainfont{Latin Modern Roman}
\usepackage{amsmath,amssymb,mathtools,unicode-math}
\setmathfont{Latin Modern Math}
\usepackage{enumitem,needspace,xurl,hyperref}
\setlength{\parindent}{0pt}
\setlength{\parskip}{5pt}
\setlength{\emergencystretch}{3em}
\newcommand{\N}{\mathbb{N}}
\newcommand{\Z}{\mathbb{Z}}
\newcommand{\Q}{\mathbb{Q}}
\newcommand{\R}{\mathbb{R}}
\newcommand{\C}{\mathbb{C}}
\newcommand{\F}{\mathbb{F}}
\newcommand{\A}{\mathbb{A}}
\newcommand{\PP}{\mathbb{P}}
\newcommand{\E}{\mathbb{E}}
\newcommand{\eps}{\varepsilon}
\newcommand{\dd}{\,\mathrm{d}}
\begin{document}
\section*{UnsolvedMath ID: 11000229}
\textbf{Source problem code:} catalogue entry
\subsection{Definitions and mathematical statement}
Let \(G<\operatorname{MCG}(S)\) be a subgroup of the mapping class group of the source's surface \(S\). The usual Thurston compactification is
\[
 \overline{\mathcal T(S)}=\mathcal T(S)\cup\partial\mathcal T(S),
\]
with boundary conventionally the projective measured-lamination space. Mapping classes act by changing markings and on that boundary. The surface hypotheses and all compactification conventions require source confirmation.

Problem 1.1 asks how the geometry of the object denoted \(\Gamma_G\) relates to the dynamics of the action of \(G\) on this whole compactification. The record never defines \(\Gamma_G\), its metric, or the geometric category. It must not be silently identified with a Cayley graph of \(G\), a particular surface-group extension, or another familiar construction. No finite-generation hypothesis on \(G\) is given.

The request is a research program: it does not specify which geometric property, dynamical behavior, invariant or implication is sought. The following comment says problems in Sections 2--4 come from the chapter author's joint paper with Benson Farb [FM02a]. The supplied summary mentions partial convex-cocompactness context but explicitly cannot verify the exact question; it supplies no theorem identifying the requested geometry/dynamics relationship.
\subsection{Short English statement}
Relate the source-defined geometry of \(\Gamma_G\) to a mapping-class subgroup's dynamics on Thurston's compactified Teichmüller space. The object \(\Gamma_G\), metric and desired relation are missing, so no specific extension-group or Cayley-graph theorem is substituted.
\subsection{Sources}
\begin{itemize}
\item Local UnsolvedMath registry, permanent ID \(11000229\), exact source code \texttt{catalogue entry}, Problem 1.1, PDF page 269. No individual proposer or proposed year is supplied; the title is truncated.
\item Farb, edited volume \emph{Problems on Mapping Class Groups and Related Topics}: \url{https://www.math.uchicago.edu/~farb/papers/mcgbook.pdf}. The volume has 329 explicitly typeset numbered items across chapters, repeated local numbering distinguished by PDF page.
\item Internal reference [FM02a], described as the chapter author's joint paper with Benson Farb and a source for problems in Sections 2--4. Full entry and the other author's identity are not supplied here.
\item Extraction pdf-font-structure, accessed 2026-07-29. Source evidence NEEDS\_REVIEW presents an open problem and requires a current primary-literature audit. Registry status open and classification OPEN-TRIAGE; summary calls it a research program, mentions partial convex-cocompactness work and explicitly cannot verify it.
\item Dataset attribution: Ulam AI / UnsolvedMath Contributors, CC BY 4.0. Source rights NEEDS\_REVIEW: publicly accessible author/publisher copy, redistribution terms requiring review.
\end{itemize}
\end{document}
